\documentclass[10pt, titlepage]{ltjsarticle}
\usepackage[utf8]{inputenc}

\usepackage[margin=23truemm]{geometry}
\usepackage{wrapfig}
\usepackage[dvipdfmx]{graphicx}
\usepackage[dvipdfmx]{color}
\usepackage{booktabs}

\usepackage{colortbl}
\usepackage{multirow}
\usepackage{tabularx}

\usepackage{enshuc}

\title{情報科学演習C \\ 課題2} 
\author{植村 向陽}
\studentid{09B24008}
\affiliation{基礎工学部情報科学科ソフトウェア科学コース}
\date{\today} 

\begin{document}

\maketitle 

\textcolor{blue}{青字の注意書きは，提出の際には削除すること．}

\textcolor{blue}{実行結果とその説明，動作をわかりやすく解説すること．その際，付録に載せたソースコードの行番号を適宜参照しながら説明すること．}

\textcolor{blue}{実行結果を載せるときは，実際にやったことが分かるように，ターミナルの結果やそのスクリーンショット等を示しつつ，説明すること．}

\section{課題2-1}
\subsection{実装方針・アルゴリズム}

\subsection{実行結果}

\subsection{考察}


\section{課題2-2}
\subsection{仕様}

\subsection{実装方針・アルゴリズム}

\subsection{実行結果}

\subsection{考察}


\section{発展課題}

\section{感想，意見，疑問等}


\appendix
\section{プログラム}
\subsection{課題2-1}
\subsection{課題2-2}
\subsection{発展課題}
\end{document}
